\scanpage{27}
Or pour $u$ variable, $\Delta(u)$ est un polynôme dans $\unsure{\mathbb Z|\mathbb Q}_p[T]$ ($T=(T_1\cdots T_{\unsure{N|\alpha}})$, $N=\deleted{\dim}\operatorname{rg}_{\mathbb Z}\operatorname{End}(A)_{\unsure{\mathrm{herm}|\mathrm{hom}}}$), qui, on vient de le voir, prend des valeurs entières pour des valeurs entières des variables. Donc $\Delta(u)$ est un polynôme à coeff. \deleted{entiers} dans $\mathbb Q$ (et non seulement dans $\mathbb Q_p$), \deleted{\unsure{donc même à coeff.} dans $\mathbb Z_{(p\mathbb Z)}$ (car $\mathbb Z_p[T]\cap\mathbb Q[T]=\mathbb Z[T]$, $\mathbb Z_p\cap\mathbb Q=\mathbb Z_{(p\mathbb Z)}$).}
Il faut prouver que ses valeurs sur le cône des $u>0$ \deleted{sont} $\geq0$, et \deleted{\unsure{pour que}} il suffit de le prouver quand on passe de $\mathbb Q$ à $\mathbb R$. Mais alors tt élém $\geq0$ est un carré d'un élém hermitien \inserted{$\unsure{v|w}$} $\geq0$, donc
\[
\Delta(u)=\Delta(\unsure{v|w}^2)=\unsure{\Delta|\delta}(\unsure{v|w})^2\in\mathbb R^2=\mathbb R^+,
\]
cqfd.

b) $\alpha_p(\varphi)$ induit bien un \uline{isom} de $\det T_p(A)_{\text{ét}}$ avec $\det T_p(A^*)_{\text{ét}}$. Cela \deleted{\unsure{exprime que ces deux}}
\inserted{équivaut à la relation}
\deleted{(i) $\alpha_p(\varphi)$ applique $\det T_p(A)_{\text{ét}}$ sur $\det T_p(A^*)_{\text{ét}}$, $\beta_p(\varphi)$}
\[
v_p\bigl(\operatorname{card}(\operatorname{Coker}\deleted{\unsure{\varphi}}T_p(\widetilde\varphi)_{\text{ét}})\bigr)=v_p\bigl(\sqrt{\deg\widetilde\varphi}\bigr)
\]
Or on a évidemment.

\mspage{26}
